\documentclass[11pt]{article}

\usepackage[margin=1in]{geometry}
\usepackage{booktabs}
\usepackage{graphicx}
\usepackage{amsmath}
\usepackage{natbib}
\usepackage{hyperref}

\title{No-Warmup Prediction-Driven Sensor Scheduling for Antarctic Blowing-Snow Monitoring}
\author{Yongzhe Li \and Zhuyu Zhang}
\date{}

\begin{document}
\maketitle

\begin{abstract}
Power-constrained automatic weather stations must decide which sensors to
operate when observing Antarctic blowing-snow processes. This draft studies a
no-warmup scheduling regime in which the strongest static allocations are less
dominant than in warm-up-constrained settings. We evaluate a prediction-driven
policy-gradient scheduler against validation-selected static allocations,
dynamic heuristics, and operationally constrained heuristic baselines. Current
evidence supports a narrow claim: removing warm-up costs can expose settings in
which learned scheduling improves over selected static allocations. Evidence
for broad dominance over dynamic heuristics remains incomplete and is treated
as an explicit claim gate rather than an assumed result.
\end{abstract}

\section{Introduction}

Antarctic blowing-snow monitoring requires multiple sensors with different
power costs and information roles. A central question is whether adaptive
scheduling improves future prediction quality when the station cannot power all
sensors simultaneously.

This manuscript focuses on the no-warmup regime. The goal is not to reuse the
existing warm-up paper narrative, but to test whether removing warm-up costs
breaks the compact-static shortcut and creates a reproducible advantage for a
prediction-driven scheduler.

\paragraph{Current claim boundary.}
The current evidence supports static-baseline advantage in many no-warmup
runs. It does not yet support a general claim of dominance over all dynamic
heuristics. The final manuscript will freeze claims only after the remaining
server-side runs are aggregated.

\section{System and Problem Formulation}

We model an automatic weather station with a set of sensors
$\mathcal{S}=\{1,\ldots,N\}$. At each time step the scheduler selects a subset
$a_t \subseteq \mathcal{S}$ subject to an instantaneous power budget. The
objective is future prediction quality over microclimate and blowing-snow
targets, not instantaneous reconstruction alone.

\section{No-Warmup Scheduling Regime}

The no-warmup setting removes startup warm-up penalties from the primary
evaluation. This makes the setting distinct from warm-up-constrained deployment
studies. It is useful as a controlled regime for testing whether static
allocations remain dominant when temporal adaptation is less penalized.

\section{Scheduler}

The learned scheduler is a prediction-driven PPO policy evaluated through the
existing PD-PPO experiment backend. The policy is compared against:

\begin{itemize}
  \item validation-selected static allocation;
  \item feasible static projected allocation;
  \item round-robin, age-of-information, and random dynamic heuristics;
  \item duty-constrained versions of the heuristic baselines.
\end{itemize}

\section{Experimental Protocol}

All training and evaluation runs are executed on the GPU server. Local scripts
only aggregate CSV/JSON outputs and generate manuscript assets.

\subsection{Evidence Gates}

We separate four gates:

\begin{enumerate}
  \item advantage over selected/static allocation;
  \item comparison to unconstrained dynamic heuristics;
  \item comparison to operationally constrained heuristics;
  \item dynamic-duty validity.
\end{enumerate}

\section{Results}

\subsection{Base No-Warmup Grid}

The partial base grid shows a strong selected/static advantage signal. Current
aggregates indicate B=1.65 wins selected/static in 9/10 seeds and B=1.70 wins
selected/static in 5/5 completed seeds. However, the same runs often lose to
the best fair dynamic heuristic and frequently fail the dynamic-duty target.

\subsection{Hard-Duty Reduced Probe}

Hard-duty constraints fix the behavioral collapse in the completed reduced
probe: seeds 41 and 42 both avoid always-on and always-off sensors. Performance
is mixed: seed41 beats selected/static but loses to round-robin and the best
constrained baseline; seed42 loses to selected/static and round-robin while
beating the best constrained baseline.

\subsection{Claim-Gate Summary}

This section will be generated from
\texttt{results/tables/claim\_gate\_summary.md} after the remaining runs are
synced and aggregated.

\section{Discussion}

The no-warmup regime weakens the strongest static shortcut, but the current
learned scheduler does not yet show robust dominance over dynamic heuristics.
This distinction is central to the final paper framing.

\section{Limitations}

The results are simulation/protocol evidence, not field deployment validation.
The no-warmup regime is a controlled scenario and should not be conflated with
warm-up-constrained physical deployment.

\section{Conclusion}

No-warmup scheduling provides a useful regime for studying when adaptive
prediction-driven scheduling can outperform selected static allocations. The
remaining experimental work determines whether the final paper can claim
operational heuristic advantage or must remain a narrower static-shortcut-break
study.

\bibliographystyle{plainnat}
\bibliography{references}

\end{document}
